\section{Analysis \& Action Plan}
\label{sec:analysis}

\paragraph{Against the definition of success.}
The routing engine meets the goal the proposal set for it. Requests return
loops close to the requested distance (US1), with real alternatives (US2), and
with a quarter to a third of the interruptions a distance-only planner would pick
(US3), at a distance cost well inside the limit. The surface preference works
reliably (US5), GPX output is valid (US7), the system is fast once warm
(US10), and the same pipeline now serves six more metro areas (US11). The
first validation run found the result weaker at the point where the runner
meets the engine. The proposal promised that a runner could ``compare
measurable tradeoffs'' and ``make the final choice'' themselves (US4), and on
16 November the interface did not show most of the tradeoffs the engine was
making. We fixed that display on 18 November and the rerun passed. Of the
seven must-haves, six are now met and US6 is met only through the served
path. One should-have is incomplete: US8 waits on the team checklist.

\paragraph{Priorities.}
We ranked issues by how much they affect what a runner sees and by whether
they break a must-have criterion. The two display gaps ranked first, so we
fixed them first. They were the only failures a runner would notice on every
request, and they undid the project's stated difference from existing tools,
which is that the tradeoffs are visible. The weak elevation preference ranks
next. It is a must-have, and
a runner who asks for hills and gets the same route as ``flat'' will conclude
the setting does nothing. Deployment issues follow, because several fixes are
already written but not live. Map-data and scaling issues rank lower. They are
real, but they mostly affect the metros beyond our primary area.

\paragraph{Decisions.}
Table~\ref{tab:actions} records a decision for each important finding. Owners
follow the areas of responsibility in our proposal: the interface for
Adithya, data, performance and deployment for Krish, and routing and scoring
for Rohan. Each item names how we will check that it worked. In every case
the check reruns the same script that found the problem, so a fix is not
counted until the measurement that failed passes.
\todo{Team to confirm owners and target dates.}

\begin{table}[ht]
\centering\footnotesize
\caption{Decisions on the main findings.}
\label{tab:actions}
\begin{tabularx}{\linewidth}{@{}>{\raggedright\arraybackslash}p{2.6cm} >{\raggedright\arraybackslash}p{1.5cm} Y l >{\raggedright\arraybackslash}p{3.1cm}@{}}
\toprule
Finding & Decision & Action & Owner & How we will verify \\
\midrule
US4: crossings, repetition, distance error and score not shown & Completed fix (\#29) & Restored crossings; added distance off target, repetition and score to each card and an expandable score breakdown & Rohan & Deployed rerun 18 Nov: 3,969 values, 0 mismatches, 8/8 components shown \\
US9: drive cost and score not shown & Completed fix (\#29) & Show driving distance and score on alternative-start cards and in the panel & Rohan & Ten-case rerun 18 Nov: both shown, 19/19 in radius \\
US6: hilly often ties flat & Design change & Let elevation influence path search, not only ranking, then retune its weight & Rohan & \texttt{acceptance.py}: US6 $\geq$80\% live \emph{and} US3 still met \\
Production backend behind \texttt{main}; manual deploys & Planned fix & Merge the reproducible release script (\#18) and redeploy from a recorded commit & Rohan & \texttt{/api/version} equals the release commit \\
Precomputed routes miss 8.7\% of requests & Completed fix, not deployed & Retry from the connected network and a wider ring (\#15) & Rohan & 600-request check $\geq$97\% after deploy \\
Live answers can be all slivers or repeat 23\% & Planned fix & Exclude routes over the repetition limit when alternatives exist; otherwise show a notice & Rohan & Shape audit; the failing park case \\
Place list served Brotli to non-browser clients & Planned fix & Stop the cache serving one encoding to every client & Krish & Both rehearsal scripts pass against the public site \\
Start locations written to logs & Planned fix & Log the area name instead of coordinates or cells & Rohan & Test that fails if a log line contains coordinates \\
Tests did not catch removed UI fields; no CI on PRs & Completed fix & Tests run on every pull request (\#21); browser tests now assert the US4 and US9 fields are present (\#29) & Rohan & CI blocks a PR that removes a required field \\
First request in a new city up to 23\,s & Backlog & Compact tile format (\#22) and pre-loading the map when a start is chosen & Krish & Deployed first-request times \\
Metros use graphs without crosswalk signals & Deferred & Rebuild with the signal-preserving builder in the next cluster batch & Krish & Signal-node counts and benchmark per metro \\
US1 5\% rule vs.\ the $\pm$0.25 mile band & Explicit decision & Keep the band; the 3 mile misses are deliberate trades for fewer interruptions & Rohan & Report both measures \\
US2 overlap definition differs from the code & Backlog & Align the duplicate test with the shorter-route measure & Rohan & US2 served rerun \\
Drive distance can be below straight line & Backlog & Add the legs from each start onto the road network & Rohan & US9 check: road $\geq$ straight line \\
\bottomrule
\end{tabularx}
\end{table}

\paragraph{Limits of the evidence.}
Several limits qualify every conclusion above.
\begin{itemize}
\item \textbf{One city.} The benchmark is 32 requests from eight starts in a
  single city. It was fixed before the system was tuned, but it is small, and
  the metro checks used six warm requests per area, where the 95th percentile
  is just the slowest request.
\item \textbf{Mapped interruptions only.} Signals and crossings are counted
  from OpenStreetMap \citep{osmwiki2026signals}. The crossing audit shows the
  planner counts what the map contains, not that the map is complete. Fewer
  signals is also not the same as safer, since a signal may protect a
  crossing.
\item \textbf{Self-built ground truth.} The crossing labels came from our own
  geometric labeller. We checked it against hand-labelled cases and inspected
  every disagreement, but it is not an independent survey.
\item \textbf{Elevation resolution.} Elevation comes from a 90\,m terrain
  model \citep{copernicus2022dem}, which captures hills but not short, steep
  grades.
\item \textbf{Test conditions.} Production timings come from one network
  path at one time of day, with no concurrent users.
\item \textbf{No outside users.} No runner outside the team has used
  RunLine under observation. Whether the tradeoffs, once displayed, actually
  help someone choose is untested, and so is whether our ranking matches
  what runners prefer.
\end{itemize}
